\section{Nonclaims and caveats}\label{sec:nonclaims}

This chapter collects the scope limits of \FIV.  The chapters state the
local caveats where they matter; here the scope of the catalog, and
answers to common questions about it, are gathered in one place.

\begin{enumerate}[label=\textbf{N\arabic*.}, leftmargin=*]
\item \textbf{\FIV\ does not re-prove substrate theorems.}
The paper states layer-agnostic laws; it does not re-derive the
carrier-specific theorems of the substrate papers.  Where a law is
anchored to or motivated by a substrate paper, the body says so, and
\Cref{tab:substrate-instances} collects these cases.

\item \textbf{\FIV\ is not a foundational reset.}
The apparatus inherited from \FoundI, \FoundII\ and \FoundIII\ is used
as background after the review in \Cref{sec:apparatus}.  Closure
formation, admissibility and quotient packaging are inputs here, not
new foundations; the \(\BirdInt\) judgment is used by no hypothesis or
proof.

\item \textbf{\FIV\ is not a unification claim.}
The axes are an organizing device for recurring obstruction types.  The
paper does not assert that they form a categorical, fibrational,
monadic or otherwise privileged structure.

\item \textbf{\FIV\ is not a complete catalog.}
The \NumLaws\ laws are the scope of this paper, not an exhaustive list.
Candidates whose recognition source, in the sense of
\Cref{sec:sketches}, has not been identified appear in
\Cref{sec:sketches}.

\item \textbf{The conditional law is conditional.}
The Hiddenness Normal Form (\Fref{F13a}) assumes its source hypothesis
\((\mathsf S)\), which is not verified here for any interface derived
from the saturated-SAT construction that motivates it.

\item \textbf{\FIV\ does not establish the substrate endpoints.}
The catalog does not prove the self-dual trace confinement hypothesis,
the Riemann hypothesis, any separation of complexity classes, or the
hiddenness results of its substrate papers
\citep{Tsiokos2026SDTC, Tsiokos2026RH, Tsiokos2026PvNP, Tsiokos2026CSL};
it states layer-agnostic laws whose instances those papers study.

\item \textbf{The catalog discipline is the contribution.}
The contribution is the uniform treatment of \NumLaws\ layer-agnostic
laws, each with its hypotheses, proof, provenance and nonclaims.  The
substrate proofs live in the cited substrate papers.

\item \textbf{The eight-axis scheme is not minimal or canonical.}
The scheme is an editorial choice, not a theorem.  Other groupings are
legitimate if they preserve the same statements and exclusions.
\end{enumerate}

The same scope answers the common questions about the catalog.

\begin{enumerate}[label=\textbf{O\arabic*.}, leftmargin=*]
\item \textbf{Why call this ``Foundations IV'' if it is a catalog?}
The title names the position in the series, not a reset.  The
contribution is the uniform statement and proof of \NumLaws\
layer-agnostic laws.

\item \textbf{Are the anchored laws merely restating substrate theorems?}
\Fref{F6} is proved here, over the reals, for any block data
satisfying its identities and budgets; its real substrate instances
use Moore--Penrose pseudoinverses of positive semidefinite forms, which
satisfy these identities.  For
\Fref{F53}, the retained saturation/zero-flag equivalence is proved
here; the source-accounting equivalence of the audited-realisability
paper is imported at its stated scope, and the counterexample of clause
(b) of \Cref{prop:F53}, a sound inference-closed ledger leaving the
obligation \(5\le2\) undischarged, is checked here, together with the
general no-go theorems and finite models.

\item \textbf{Why is the Hiddenness Normal Form stated conditionally?}
Its first four conjuncts concern declared predicates of an interface
whose intended source is the saturated-SAT construction of
\citet{Tsiokos2026PvNP}.  No such interface is constructed there or
here, so these conjuncts enter as the hypothesis \((\mathsf S)\); the
fifth conjunct and the equivalence of surplus with non-collapse are
proved.

\item \textbf{Is the hypothesis \((\mathsf S)\) consistent?}
Yes: it holds for small explicit interfaces, including one with a
predictive surplus (\Cref{sec:access:hiddenness}).  This shows that
part (i) of \Cref{thm:F13a} is not vacuous; it says nothing about the
intended saturated-SAT interface.

\item \textbf{Why a single paper rather than one per axis?}
A single paper states the apparatus once and lets laws of different
axes be compared in one register.

\item \textbf{Why these \NumLaws\ laws and not others?}
They are the recurring structural obstructions and closures of the
Six Birds series.  The catalog is not exhaustive; \Cref{sec:sketches}
records candidates that remain outside it.

\item \textbf{Why do the Status/Records/Coherence and
Self-Reference/Limits chapters hold so many laws?}
Those axes cover the broadest range of record, coherence, information,
limit and self-reference phenomena.  Short per-law subsections and the
summary tables keep them navigable.

\item \textbf{Do the apparatus restatements lose nuance from the
foundations papers?}
Yes; the review in \Cref{sec:apparatus} is intentionally compact.  It
suffices for reading this paper, and the cited foundations papers carry
the operational detail.
\end{enumerate}
